%% abtex2-modelo-trabalho-academico.tex, v-1.9.6 laurocesar
%% Copyright 2012-2016 by abnTeX2 group at http://www.abntex.net.br/ 
%%
%% This work may be distributed and/or modified under the
%% conditions of the LaTeX Project Public License, either version 1.3
%% of this license or (at your option) any later version.
%% The latest version of this license is in
%%   http://www.latex-project.org/lppl.txt
%% and version 1.3 or later is part of all distributions of LaTeX
%% version 2005/12/01 or later.
%%
%% This work has the LPPL maintenance status `maintained'.
%% 
%% The Current Maintainer of this work is the abnTeX2 team, led
%% by Lauro César Araujo. Further information are available on 
%% http://www.abntex.net.br/
%%
%% This work consists of the files abntex2-modelo-trabalho-academico.tex,
%% abntex2-modelo-include-comandos and abntex2-modelo-references.bib
%%

% ------------------------------------------------------------------------
% ------------------------------------------------------------------------
% abnTeX2: Modelo de Trabalho Academico (tese de doutorado, dissertacao de
% mestrado e trabalhos monograficos em geral) em conformidade com 
% ABNT NBR 14724:2011: Informacao e documentacao - Trabalhos academicos -
% Apresentacao
% ------------------------------------------------------------------------
% ------------------------------------------------------------------------

\documentclass[
	% -- opções da classe memoir --
	12pt,				% tamanho da fonte
	openright,			% capítulos começam em pág ímpar (insere página vazia caso preciso)
	twoside,			% para impressão em recto e verso. Oposto a oneside
	a4paper,			% tamanho do papel. 
	% -- opções da classe abntex2 --
	chapter=TITLE,		% títulos de capítulos convertidos em letras maiúsculas
	section=TITLE,		% títulos de seções convertidos em letras maiúsculas
	%subsection=TITLE,	% títulos de subseções convertidos em letras maiúsculas
	%subsubsection=TITLE,% títulos de subsubseções convertidos em letras maiúsculas
	% -- opções do pacote babel --
	english,			% idioma adicional para hifenização
	french,				% idioma adicional para hifenização
	spanish,			% idioma adicional para hifenização
	brazil				% o último idioma é o principal do documento
	]{abntex2}

% ---
% Pacotes básicos 
% ---
\usepackage{lmodern}			% Usa a fonte Latin Modern			
\usepackage[T1]{fontenc}		% Selecao de codigos de fonte.
\usepackage[utf8]{inputenc}		% Codificacao do documento (conversão automática dos acentos)
\usepackage{lastpage}			% Usado pela Ficha catalográfica
\usepackage{indentfirst}		% Indenta o primeiro parágrafo de cada seção.
\usepackage{color}				% Controle das cores
\usepackage{graphicx}			% Inclusão de gráficos
\usepackage{microtype} 			% para melhorias de justificação
\usepackage{subfiles}
\usepackage{longtable}
\usepackage{hyperref}
\usepackage{setspace}
\usepackage{hyphenat}
\usepackage{float}
\usepackage{longtable}
\usepackage{caption}
\usepackage{subcaption}
\usepackage{gensymb}

 
%\hyphenation{PROCFIS, diâmetro, condu-tân-cia, cor-ren-te, assoa-lho, noradre-nalina, den-tea-do}
% ---
		
% ---
% Pacotes adicionais, usados apenas no âmbito do Modelo Canônico do abnteX2
% ---
\usepackage{lipsum}				% para geração de dummy text
% ---

% ---
% Pacotes de citações
% ---
\usepackage[brazilian,hyperpageref]{backref}	 % Paginas com as citações na bibl
\usepackage[alf]{abntex2cite}	% Citações padrão ABNT

% --- 
% CONFIGURAÇÕES DE PACOTES
% --- 

% ---
% Configurações do pacote backref
% Usado sem a opção hyperpageref de backref
\renewcommand{\backrefpagesname}{Citado na(s) página(s):~}
% Texto padrão antes do número das páginas
\renewcommand{\backref}{}
% Define os textos da citação
\renewcommand*{\backrefalt}[4]{
	\ifcase #1 %
		Nenhuma citação no texto.%
	\or
		Citado na página #2.%
	\else
		Citado #1 vezes nas páginas #2.%
	\fi}%
% ---
% \renewcommand{\folhaderostocontent}{
% \begin{center}
% {\includegraphics[scale=.2]{figuras/ufs_vertical_positiva} \\
% \ABNTEXchapterfont\large\imprimirautor}
% \vspace*{\fill}\vspace*{\fill}
% \begin{center}
% \ABNTEXchapterfont\bfseries\Large\imprimirtitulo
% \end{center}
% \vspace*{\fill}
% \abntex@ifnotempty{\imprimirpreambulo}{%
% \hspace{.45\textwidth}
% \begin{minipage}{.5\textwidth}
% \SingleSpacing
% \imprimirpreambulo
% \end{minipage}%
% \vspace*{\fill}
% }%
% {\abntex@ifnotempty{\imprimirinstituicao}{\imprimirinstituicao
% \vspace*{\fill}}}
% {\large\imprimirorientadorRotulo~\imprimirorientador\par}
% \abntex@ifnotempty{\imprimircoorientador}{%
% {\large\imprimircoorientadorRotulo~\imprimircoorientador}%
% }%
% \vspace*{\fill}
% {\large\imprimirlocal}
% \par
% {\large\imprimirdata}
% \vspace*{1cm}
% \end{center}
% }
% \makeatother



% ---
% Informações de dados para CAPA e FOLHA DE ROSTO
% ---
\titulo{Título da dissertação}
\autor{Nome do aluno }
\local{S\~ao Cristóvão  - SE }
\data{2021}
\orientador{Prof. Dr. Nome do professor}
% \coorientador{Equipe \abnTeX}
\instituicao{%
  Universidade Federal de Sergipe
  \par
  Departamento de XXXXXXX
  \par
  Programa de Pós-Graduação em XXXXXXXX 
  \par
  Sigla do programa
  }
\tipotrabalho{Tese (Doutorado)}
% O preambulo deve conter o tipo do trabalho, o objetivo, 
% o nome da instituição e a área de concentração 
\preambulo{Tese apresentada à Banca de Defesa do Programa de Pós-Graduação em XXXXX – Sigla, como parte dos requisitos para aspirar ao título de Doutor em Ciências XXXXXX. }
% ---


% ---
% Configurações de aparência do PDF final

% alterando o aspecto da cor azul
\definecolor{blue}{RGB}{41,5,195}

% informações do PDF
\makeatletter
\hypersetup{
     	%pagebackref=true,
		pdftitle={\@title}, 
		pdfauthor={\@author},
    	pdfsubject={\imprimirpreambulo},
	    pdfcreator={LaTeX with abnTeX2},
		pdfkeywords={abnt}{latex}{abntex}{abntex2}{trabalho acadêmico}, 
		colorlinks=true,       		% false: boxed links; true: colored links
    	linkcolor=black,          	% color of internal links
    	citecolor=black,        		% color of links to bibliography
    	filecolor=magenta,      		% color of file links
		urlcolor=blue,
		bookmarksdepth=4
}
\makeatother
% --- 

% --- 
% Espaçamentos entre linhas e parágrafos 
% --- 

% O tamanho doo parágrafo é dado por:
\setlength{\parindent}{1.3cm}

% Controle do espaçamento entre um parágrafo e outro:
\setlength{\parskip}{0.2cm}  % tente também \onelineskip

% ---
% compila o indice
% ---
\makeindex
% ---

%%% -----
%%% Formato de cabeçalho/rodapé romano nos elementos pré-textuais
%%% -----

%% Novo estilo
\makepagestyle{estilo_pretextual} %%% escolha um nome
  \makeevenhead{estilo_pretextual}{}{}{\ABNTEXfontereduzida \textbf \thepage}
  \makeoddhead{estilo_pretextual}{}{}{\ABNTEXfontereduzida \textbf \thepage}

%% Customiza comando \pretextual
\renewcommand{\pretextual}{
  \pagenumbering{roman} %%% ou \pagenumbering{Roman}
  \aliaspagestyle{chapter}{estilo_pretextual}% customizing chapter pagestyle
  \pagestyle{estilo_pretextual}
  \aliaspagestyle{cleared}{empty}
  \aliaspagestyle{part}{estilo_pretextual}
}

% ---
% Ajusta a marca \textual para que a numeração volte a ser arábica
% nos elementos textuais
\let\oldtextual\textual        % copia o comando \textual anterior para \oldtextual
\renewcommand{\textual}{%
  \oldtextual%
  \pagenumbering{arabic} % volta à numeração arábica
}
% ---



% ----
% Início do documento
% ----
\begin{document}
\pagenumbering{roman}
% \chead{\thepage}
% \pagestyle{empty}

% Seleciona o idioma do documento (conforme pacotes do babel)
%\selectlanguage{english}
\selectlanguage{brazil}

% Retira espaço extra obsoleto entre as frases.
\frenchspacing 

% ----------------------------------------------------------
% ELEMENTOS PRÉ-TEXTUAIS
% ----------------------------------------------------------
% \pretextual

% ---
% Capa
% ---
\imprimircapa
% ---

% ---
% Folha de rosto
% (o * indica que haverá a ficha bibliográfica)
% ---
\imprimirfolhaderosto*
% ---

% ---
% Inserir a ficha bibliografica
% ---

% Isto é um exemplo de Ficha Catalográfica, ou ``Dados internacionais de
% catalogação-na-publicação''. Você pode utilizar este modelo como referência. 
% Porém, provavelmente a biblioteca da sua universidade lhe fornecerá um PDF
% com a ficha catalográfica definitiva após a defesa do trabalho. Quando estiver
% com o documento, salve-o como PDF no diretório do seu projeto e substitua todo
% o conteúdo de implementação deste arquivo pelo comando abaixo:
%
% \begin{fichacatalografica}
%     \includepdf{fig_ficha_catalografica.pdf}
% \end{fichacatalografica}

\begin{fichacatalografica}
	\sffamily
	\vspace*{\fill}					% Posição vertical
	\begin{center}					% Minipage Centralizado
	\fbox{\begin{minipage}[c][8cm]{15cm}		% Largura
	\small
	\imprimirautor
	%Sobrenome, Nome do autor
	
	\hspace{0.5cm} \imprimirtitulo  / \imprimirautor. --
	\imprimirlocal, \imprimirdata-
	
	\hspace{0.5cm} \pageref{LastPage} p. : il. (algumas color.) ; 30 cm.\\
	
	\hspace{0.5cm} \imprimirorientadorRotulo~\imprimirorientador\\
	
	\hspace{0.5cm}
	\parbox[t]{\textwidth}{\imprimirtipotrabalho~--~\imprimirinstituicao,
	\imprimirdata.}\\
	
	\hspace{0.5cm}
		1. Palavra chave.
		2. Palavra chave.
		3. Palavra chave.
		I. Orientador: \imprimirorientador \hspace{0.1mm} 
		II. Universidade Federal de Sergipe.
		III. Programa de Pós-Graduação em XXXXX - Sigla.
		IV.  \imprimirtitulo		
	\end{minipage}}
	\end{center}
\end{fichacatalografica}
% ---
% 
% % ---
% % Inserir errata
% % ---
% \begin{errata}
% Elemento opcional da \citeonline[4.2.1.2]{NBR14724:2011}. Exemplo:
% 
% \vspace{\onelineskip}
% 
% FERRIGNO, C. R. A. \textbf{Tratamento de neoplasias ósseas apendiculares com
% reimplantação de enxerto ósseo autólogo autoclavado associado ao plasma
% rico em plaquetas}: estudo crítico na cirurgia de preservação de membro em
% cães. 2011. 128 f. Tese (Livre-Docência) - Faculdade de Medicina Veterinária e
% Zootecnia, Universidade de São Paulo, São Paulo, 2011.
% 
% \begin{table}[htb]
% \center
% \footnotesize
% \begin{tabular}{|p{1.4cm}|p{1cm}|p{3cm}|p{3cm}|}
%   \hline
%    \textbf{Folha} & \textbf{Linha}  & \textbf{Onde se lê}  & \textbf{Leia-se}  \\
%     \hline
%     1 & 10 & auto-conclavo & autoconclavo\\
%    \hline
% \end{tabular}
% \end{table}
% 
% \end{errata}
% ---

% ---
% Inserir folha de aprovação
% ---

% Isto é um exemplo de Folha de aprovação, elemento obrigatório da NBR
% 14724/2011 (seção 4.2.1.3). Você pode utilizar este modelo até a aprovação
% do trabalho. Após isso, substitua todo o conteúdo deste arquivo por uma
% imagem da página assinada pela banca com o comando abaixo:
%
% \includepdf{folhadeaprovacao_final.pdf}
%
\begin{folhadeaprovacao}

  \begin{center}
    {\ABNTEXchapterfont\large\imprimirautor}

    \vspace*{\fill}\vspace*{\fill}
    \begin{center}
      \ABNTEXchapterfont\bfseries\Large\imprimirtitulo
    \end{center}
    \vspace*{\fill}
    
    \hspace{.45\textwidth}
    \begin{minipage}{.5\textwidth}
        \imprimirpreambulo
    \end{minipage}%
    \vspace*{\fill}
   \end{center}
        
%    Trabalho que será apresentado. \imprimirlocal, xx de xxxxx de 2018:

   \assinatura{\textbf{Professor} \\ Convidado 1} 
   \assinatura{\textbf{Professor} \\ Convidado 2}
   \assinatura{\textbf{Professor} \\ Convidado 3}
   \assinatura{\textbf{Professor} \\ Convidado 3}
   \assinatura{\textbf{Professor} \\ Convidado 4}
      
   \begin{center}
    \vspace*{0.5cm}
    {\large\imprimirlocal}
    \par
    {\large\imprimirdata}
    \vspace*{1cm}
  \end{center}
  
\end{folhadeaprovacao}
% ---

% ---
% Dedicatória
% ---
\begin{dedicatoria}
   \vspace*{\fill}
   \centering
   \noindent
   \textit{ } \vspace*{\fill}
\end{dedicatoria}
% ---

% ---
% Agradecimentos
% ---
\begin{agradecimentos}



\end{agradecimentos}
% ---

% ---
% Epígrafe
% ---
% \begin{epigrafe}
%     \vspace*{\fill}
% 	\begin{flushright}
% 		\textit{ }
% 	\end{flushright}
% \end{epigrafe}
% ---

% ---
% RESUMOS
% ---

% resumo em português
\setlength{\absparsep}{8pt}
\linespread{1.0}
% ajusta o espaçamento dos parágrafos do resumo
\begin{resumo}[Resumo]

{\footnotesize
Resumo xxx xxx xxx xxxxxx xxxx
}
 \vspace{\onelineskip}
 
  \textbf{Palavras-chave}: Palavra-chave 1, palavra-chave 2.

\end{resumo}

% resumo em inglês
\linespread{1.0}
\begin{resumo}[Abstract]
 \begin{otherlanguage*}{english}
 {\footnotesize
Abstract xxxx xxxx 

}
   \vspace{\onelineskip}
 
   \noindent 
   \textbf{Keywords}: keyword 1, keyword 2.
 \end{otherlanguage*}
\end{resumo}


% ---

% ---
% inserir lista de ilustrações
% ---
\pdfbookmark[0]{\listfigurename}{lof}
\listoffigures*
\cleardoublepage
% ---

% ---
% inserir lista de tabelas
% ---
\pdfbookmark[0]{\listtablename}{lot}
\listoftables*
\cleardoublepage
% ---

% ---
% inserir lista de abreviaturas e siglas
% ---
\begin{siglas}
  \item
  
\end{siglas}
% ---

% ---
% inserir lista de símbolos
% ---
\begin{simbolos}

  \item
%   \item[
\end{simbolos}
% ---

% ---
% inserir o sumario
% ---
\pdfbookmark[0]{\contentsname}{toc}
\tableofcontents*
\cleardoublepage
% ---



% ----------------------------------------------------------
% ELEMENTOS TEXTUAIS
% ----------------------------------------------------------
\textual
\pagenumbering{arabic}
% ----------------------------------------------------------
\subfile{Introducao}
% ----------------------------------------------------------

\subfile{Revisao}

%  \chapter*[Introdução]{Introdução}
% \addcontentsline{toc}{chapter}{Introdução}
% ---------------------------------------------------------- 
% % Este documento e seu código-fonte são exemplos de referência de uso da classe
%  \textsf{abntex2} e do pacote \textsf{abntex2cite}. O documento 
% % exemplifica a elaboração de trabalho acadêmico (tese, dissertação e outros do
% gênero) produzido conforme a ABNT NBR 14724:2011 \emph{Informação e documentação
% - Trabalhos acadêmicos - Apresentação}.
\subfile{Objetivos}

\subfile{Material}
% 
\subfile{Resultados}
% \cite{abntex2classe,abntex2cite,abntex2cite-alf} e da classe \textsf{memoir}
% \cite{memoir}. 
%
% Esperamos, sinceramente, que o \abnTeX\ aprimore a qualidade do trabalho que
% você produzirá, de modo que o principal esforço seja concentrado no principal:
% na contribuição científica.
 
% Equipe \abnTeX 
% 
% Lauro César Araujo
% 
% % ----------------------------------------------------------
% % PARTE
% % -----------------------------------------------------
\subfile{Discussao}
% % ----------------------------------------------------------
% 
% % ---
% % Capitulo com exemplos de comandos inseridos de arquivo externo 
% % ---
% \include{abntex2-modelo-include-comandos}
% % ---
% 
% % ----------------------------------------------------------
% % PARTE
% % ----------------------------------------------------------
% \part{Referenciais teóricos}
% % ----------------------------------------------------------
% 

% % ---
% 
% \lipsum[1]
% 
% \lipsum[2-3]
% 
% % ----------------------------------------------------------
% % PARTE
% % ----------------------------------------------------------
% \part{Resultados}
% % ----------------------------------------------------------
% 
% % ---
% % primeiro capitulo de Resultados
% % ---
% \chapter{Lectus lobortis condimentum}
% % ---
% 
% % ---
% \section{Vestibulum ante ipsum primis in faucibus orci luctus et ultrices
% posuere cubilia Curae}
% % ---
% 
% \lipsum[21-22]
% 
% % ---
% % segundo capitulo de Resultados
% % ---
% \chapter{Nam sed tellus sit amet lectus urna ullamcorper tristique interdum
% elementum}
% % ---
% 
% % ---
% \section{Pellentesque sit amet pede ac sem eleifend consectetuer}
% % ---
% 
% \lipsum[24]
% 
% % ----------------------------------------------------------
% % Finaliza a parte no bookmark do PDF
% % para que se inicie o bookmark na raiz
% % e adiciona espaço de parte no Sumário
% % ----------------------------------------------------------
% \phantompart
% 
% % ---
% % Conclusão
% % ---
\subfile{Conclusao}
% % ---
% 
% \lipsum[31-33]
% 
% % ----------------------------------------------------------
% % ELEMENTOS PÓS-TEXTUAIS
% % ----------------------------------------------------------
% \postextual
% ----------------------------------------------------------

% ----------------------------------------------------------
% Referências bibliográficas
% ----------------------------------------------------------
\bibliography{Bibliografia.bib}

% ----------------------------------------------------------
% Glossário
% ----------------------------------------------------------
%
% Consulte o manual da classe abntex2 para orientações sobre o glossário.
%
%\glossary

% ----------------------------------------------------------
% Apêndices
% ----------------------------------------------------------

% ---
% Inicia os apêndices
% ---
\begin{apendicesenv}
\partapendices
\subfile{Anexos1}

% 
% % Imprime uma página indicando o início dos apêndices

% 
% % ----------------------------------------------------------
% \chapter{Primeiro apêndice}
% % ----------------------------------------------------------
% 
% \lipsum[50]
% 
% % ----------------------------------------------------------
% % \chapter{Nullam elementum urna vel imperdiet sodales elit ipsum pharetra ligula
% % ac pretium ante justo a nulla curabitur tristique arcu eu metus}
% % ----------------------------------------------------------
% \lipsum[55-57]
% 
\end{apendicesenv}
% % ---
% 
% 
% % ----------------------------------------------------------
% % Anexos
% % ----------------------------------------------------------

% ---
% Inicia os anexos
% ---
% \begin{anexosenv}

% Imprime uma página indicando o início dos anexos
% \partanexos
% \subfile{Anexos1}
% \subfile{Anexos2}
% \subfile{Anexos3}
% \subfile{Anexos4}
% \subfile{Anexos5}
% \subfile{Anexos6}
% \subfile{Anexos7}
% \subfile{Anexos8}
% \subfile{Anexos9}
% \subfile{tabela}

% % ---
% \chapter{Morbi ultrices rutrum lorem.}
% % ---
% \lipsum[30]
% 
% % ---
% \chapter{Cras non urna sed feugiat cum sociis natoque penatibus et magnis dis
% parturient montes nascetur ridiculus mus}
% % ---
% 
% \lipsum[31]
% 
% % ---
% \chapter{Fusce facilisis lacinia dui}
% % ---
% 
% \lipsum[32]
% 
% \end{anexosenv}

%---------------------------------------------------------------------
% INDICE REMISSIVO
%---------------------------------------------------------------------
\phantompart
\printindex
%---------------------------------------------------------------------

\end{document}
